%=== divers/PR/loiproba_01.tex || PR || Loi de proba

On lance un dé truqué à $6$ faces. Les probabilités sont données dans le tableau suivant :

\smallskip

\begin{tblr}{width=0.975\linewidth,hlines,vlines,colspec={X[m,c]X[m,c]X[m,c]X[m,c]X[m,c]X[m,c]}}
  Face n°$1$ & Face n°$2$ & Face n°$3$ & Face n°$4$ & Face n°$5$ & Face n°$6$ \\
  $\dfrac{1}{12}$ & $\dfrac{1}{4}$ & $\dfrac{1}{12}$ & $\dfrac{1}{6}$ & $\dfrac{1}{12}$ & $x$ \\
\end{tblr}

\smallskip

On peut affirmer que :

\smallskip
\eamreponsesautom[2]{$x=\dfrac{1}{4}$}{$x=\dfrac{1}{12}$}{$x=\dfrac{1}{3}$}{$x=\dfrac{1}{6}$}
